\section{Introduction}\label{sec:intro}

Each element added to a population-dynamics model is a claim that the structure it encodes helps to predict the system. In stock assessment such claims are usually examined through fit, retrospective pattern and simulation self-consistency \citep{HilbornWalters1992,Punt2016,Maunder2015}; where prediction skill is examined directly, through hindcast cross-validation, the quantity predicted is the fitted survey or catch-rate index \citep{Kell2016,Kell2021,Carvalho2021,Kokkalis2024}. Advice, however, concerns the state: spawning-stock biomass (SSB) and its position relative to a reference point. The forecasting literature supplies the missing benchmark. A model is held to earn its complexity only if it beats a naive forecast, typically last-value persistence, out of sample at the horizons that matter \citep{Makridakis2000,HyndmanKoehler2006,HyndmanAthanasopoulos2021}; evaluation is run on a rolling origin so that each forecast uses only information available at the time it would have been issued \citep{Tashman2000,Bergmeir2012}; and differences in accuracy are tested with the Diebold--Mariano statistic and its small-sample corrections \citep{Diebold1995,Harvey1997,Diebold2015}. In ecological forecasting these standards have produced sobering results: simple time-series models forecast fish populations about as well as complex ones \citep{Ward2014}, and model-free methods have outperformed the mechanistic models that generated the data when the mechanism is unstable \citep{Perretti2013}. Whether a structural module improves state forecasts for a particular stock is therefore an empirical question, and one that can be settled by a scored test rather than argued.

This paper reports such a test on Northern cod (\emph{Gadus morhua}) in NAFO Divisions 2J3KL. The stock is the archetype of a fishery collapse: SSB fell by an order of magnitude between 1990 and 1994 under catches of roughly 200\,kt, a moratorium was declared in July 1992, the causes have been debated ever since \citep{Hutchings1994,Walters1996,Myers1997,Shelton1999}, and so, more recently, have the pace and the reality of its recovery \citep{RoseRowe2015,RoseWalters2019,Rose2026}. It is also an archetype of the reference-point problem. The limit reference point was for many years defined as the 1983--1989 mean SSB from the Northern Cod Assessment Model \citep[NCAM;][]{Cadigan2016,DFO2016}; a revised framework in 2023--2024 extended the assessed series back to 1954 and replaced the limit with a $\Blim$ of 276\,kt \citep{DFO2024a,DFO2024b,Regular2025}, under which a stock that had been classified in the critical zone for over two decades was reclassified at once. Two retrospective SSB series therefore exist, NCAM for 1983--2015 and xteNCAM for 1954--2024, each with its own catch record and threshold, and each a legitimate state on which to score a forecast. Neither is a truth; they are the two series that advice has been built on, and a method that claims forecast skill should show it on at least one of them.

The models tested form a ladder of nested surplus-production forecasters \citep{Schaefer1954,Pella1969}: an autonomous module driven by the training-mean catch, a depensatory variant, a stock-flow module that receives the realised catch along the horizon, an autoregressive residual correction, and a lagged-start variant. The ladder is deliberately simple. The question is not whether a surplus-production model can be made to fit Northern cod (it can, and the difficulties of fitting such models are well documented; \citealp{LudwigWalters1985,Polacheck1993,Punt2003,PedersenBerg2017}) but whether any of its structural steps lowers forecast error on SSB relative to the step below it and relative to persistence, at one- and five-year horizons, under a declared rule. The rule has three parts: a comparator hurdle against the next-simpler module, a persistence hurdle, and a requirement that both hold at both horizons, with a 5\% band for ties. Its scoring core (the windows, the rolling-origin RMSE comparison and the persistence bar) was coded before the first scoring pass and applied unchanged; the tie band and the comparator declarations for the non-nested modules were formalised after the first scores had been seen, and neither decides any outcome reported here. The design is a fixed computational protocol with a recorded pass order, not a prospective registration (SI S1), and it is applied to both assessment series without pooling them.

The result is uniform: no module is retained on either series. A verdict of this kind is only as informative as the test is powerful, and a pooled score can hide where the deficit arises. The paper therefore goes beyond the verdict in three ways that are new in this revision. First, the pooled score is decomposed by period, which shows that the deficit is carried by the 1991--1995 collapse, that the autonomous modules forecast the numerical floor at every origin after 1995 (so that their post-collapse errors are simply the observed biomass), and that the catch-supplied modules have one-step skill in the low-biomass regime on one specification but not on the other. Second, the failure to forecast the collapse is traced to an identification result: the lower equilibrium of a catch-driven Schaefer map is generated by the removals, not by depensation, so no autonomous member of the class can produce a collapse from a pre-collapse training window, whatever its Allee term. Third, the operating characteristics of the rule are measured by simulation under known truth, which shows that non-retention is strong evidence against the autonomous module, weak evidence against the stock-flow and depensatory modules, and that the rule over-retains when the truth lies outside the class. Finally, six outcome years after the first series ends, together with the 2023 framework revision and the 2022--2023 survey suspension, are used to ask which layer of the advice problem moved: forecast structure, or the observation and reference-point layer that any forecast is scored against.

Section~\ref{sec:data} describes the two specifications and the information available at each forecast origin. Section~\ref{sec:methods} defines the ladder, the scores, the retention rule and the uncertainty and simulation designs. Section~\ref{sec:results} reports the pooled verdict, the fixed-window diagnostics, the period decomposition, the variant passes, the identification argument, the operating characteristics and the outcome years. Section~\ref{sec:discussion} interprets the findings and states their limits; Section~\ref{sec:conclusions} concludes. Material that supports but does not carry the argument (full test grids, fitted parameters, the floor-forecast census, secondary scores, the simulation design, the suspension chronology and the verification record) is in the Supplementary Information (SI).
